\subsubsection{Bond Future Option}
\label{ss:input_commodity_option}

\ifdefined\IncludePayoff
{
\textbf{Payoff}

A bond future option is an exchange traded option giving the holder the option to enter into an underlying bond future 
contract at a specified price called the strike price. The underlying bond future contract that is entered into is 
described in Section~\ref{ss:bondfuture}. The options may be European or American exercise and are generally 
short dated options on upcoming future contract months.

{\bf Input}
}
\fi

The \lstinline!BondFutureOptionData! node is the trade data container for the \lstinline!BondFutureOption! trade type.
Vanilla bond future options with exercise style \lstinline!European! or \lstinline!American! are supported.
The \lstinline!BondFutureOptionData! node includes exactly one \lstinline!OptionData! trade component sub-node 
along with elements specific to the bond future option.
The structure of a \lstinline!BondFutureOptionData! node for a bond future option is shown in 
Listing \ref{lst:bond-future-option_data}.
Note that the \lstinline!BondFutureOption! trade type can be used as a component in an index or basket underlying a 
Generic TRS. The Generic TRS is described in Section~\ref{ss:GenericTRS}. In particular, the 
\lstinline!BondFutureOption! node can be included as a component in the \lstinline!Derivative! node or 
as a component in the basket referenced by the \lstinline!PortfolioIndexTradeData! node.

\begin{listing}[H]
%\hrule\medskip
\begin{minted}[fontsize=\footnotesize]{xml}
<BondFutureOptionData>
  <OptionData>
    <LongShort>Short</LongShort>
    <OptionType>Call</OptionType>
    <Style>American</Style>
    <ExerciseDates>
      <ExerciseDate>2026-04-10</ExerciseDate>
    </ExerciseDates>
  </OptionData>
  <ContractName>TYM26</ContractName>
  <ContractNotional>1000000</ContractNotional>
  <Strike>1.1175</Strike>
</BondFutureOptionData>
\end{minted}
\caption{Bond Future Option data}
\label{lst:bond-future-option_data}
\end{listing}

The meaning and allowable values of the elements in the \lstinline!BondFutureOptionData! node are given below.

\begin{itemize}

\item \lstinline!OptionData!:
  The \lstinline!OptionData! node is as described in Section~\ref{ss:option_data}.
  The relevant fields in the \lstinline!OptionData! node for a \lstinline!BondFutureOption! are:
  \begin{itemize}
  \item \lstinline!LongShort!: The allowable values are \lstinline!Long! or \lstinline!Short!.
  \item \lstinline!OptionType!: The allowable values are \lstinline!Call! or \lstinline!Put!.
  \item \lstinline!Style!: The allowable values are \lstinline!European! or \lstinline!American!.
  \item \lstinline!ExerciseDates!:
    This node must contain exactly one \lstinline!ExerciseDate! element giving the option expiry date.
    Allowable values: See Date description in Table \ref{tab:allow_stand_data}.
  \item \lstinline!Premiums! [Optional]:
    Option premium node with amounts paid by the option buyer to the option seller.
    Allowable values: see Section~\ref{ss:premiums}
  \end{itemize}
 
\item \lstinline!ContractName!: The name of the underlying bond future contract. \\
  Allowable values: see \lstinline!ContractName! for bond future trades in Section~\ref{ss:bondfuture}.

\item \lstinline!ContractNotional!: The notional of the position expressed in the currency of the bond future contract. \\
  Allowable values: A non-negative real number.

\item \lstinline!Strike!: The bond future option strike price.
  Note: bond future prices and option strikes are generally expressed in points and fractions of points.
        The \lstinline!Strike! expects the price in absolute terms.
        For example, a strike price of 111.75 points is given in the \lstinline!Strike! node as \lstinline!1.1175!. \\
  Allowable values: A positive real number.

\end{itemize}

Note also that \lstinline!SeparateCallPutVols! can be set in the additional fields section of the trade envelope to 
override the setting of \lstinline!SeparateCallPutVols! given in the pricing engines configuration.
This parameter can be set to true or false or to a regex inclusion pattern.
For the trade envelope, setting it to \lstinline!true! or \lstinline!false! is probably all that is needed. In other 
words, the regex option is not that useful at the trade level and was only added to support a more granular configuration 
at the pricing engine level.
If set to false, the default, the same volatility surface is used for puts and calls.
If set to true, there will be a separate volatility surface used to price puts and calls on the same underlying 
bond future contract. Which surface is picked depends on the option type of the trade.
If set to a regex pattern, all bond future option trades where the underlying contract name matches the regex 
will use a volatility surface for puts if the option type is a put and a volatility surface for calls if the 
option type is a call.
Obviously, the bond future volatility structures need to be configured in a way that matches the setting here.
See the bond future volatility curve configuration for details on how to set these up. The convention is 
generally to use the underlying bond future contract name as the curve ID of the bond future volatility 
curve where only one surface is needed for that contract.
Where a separate call and put surface is needed for the contract, the bond future contract name should be 
suffixed with \lstinline!_CALL! and \lstinline!_PUT! respectively.
